\section{A prime-slot shifted correlation with a common parity twist}
\label{sh:section}

We prove the shifted-correlation estimate needed below.  The graph
transference and residual ideal operator are imported from
\cite{SmoothShifted}; the replacement of a long rough-integer factor by
a prime factor, and the common parity twist, are proved here.

\subsection{Groups, marked weights, and the assertion}

Fix
\[
 0<a<b<c<d<0.47,\qquad c_0,C_0,v_-,v_+>0,
 \qquad 0<q_0<1,\qquad \ell\geq1.
\]
The integer $\ell$ is fixed.  Let $\mathcal S,\mathcal B$ index
pairwise disjoint groups of primes $\mathcal P_g$, with
\begin{gather}
 c_0T\leq |\mathcal S|,|\mathcal B|\leq C_0T,
 \qquad v_-\leq V_g:=\sum_{p\in\mathcal P_g}\frac1p\leq v_+,
 \label{sh:groups}\\
 \mathcal P_g\subset[\exp(L^a),\exp(L^b)]\quad(g\in\mathcal S),
 \qquad
 \mathcal P_g\subset[\exp(L^c),\exp(L^d)]\quad(g\in\mathcal B).
 \notag
\end{gather}
Each big group is the set of all primes in an interval.  Put
$\mathcal P=\bigcup_g\mathcal P_g$, $s_P=|\mathcal S|+|\mathcal B|$,
and, for $n\geq1$,
\[
 n_*=\prod_{p\notin\mathcal P}p^{v_p(n)},\qquad
 \omega_g(n)=\sum_{p\in\mathcal P_g}\1_{p\mid n},\qquad
 \omega_P(n)=\sum_g\omega_g(n).
\]
For the graph definitions below, extend $\omega_g$ and $\omega_P$
to all $n\in\mathbb Z$ by these same divisibility formulas, including
the convention that every prime divides zero.
For a vector $\mathbf t=(t_g)$ of nonnegative integers define
\begin{equation}\label{sh:marked-weight}
 W_{\mathbf t}(n)=q_0^{\omega_P(n)-\sum_g t_g}
                   \prod_g\frac{(\omega_g(n))_{t_g}}{V_g^{t_g}},
 \qquad W_\ell=W_{(\ell,\ldots,\ell)}.
\end{equation}
Here $(u)_t=u(u-1)\cdots(u-t+1)$, and $(u)_0=1$.
An ordered marked list has $t_g$ distinct prime divisors from each
$\mathcal P_g$.  If $\mathbf b$ is a function of these lists, define
\[
 W_{\mathbf t}^{\mathbf b}(n)
 =q_0^{\omega_P(n)-\sum_g t_g}
       \prod_g V_g^{-t_g}\sum_{\text{marked lists at }n}\mathbf b.
\]
Equivalently, this is $W_{\mathbf t}(n)$ times the average of
$\mathbf b$ over its lists, with value zero when no such list exists.
If $t_g\leq t_{\max}$ and $|\mathbf b|\leq1$, then
\begin{equation}\label{sh:weight-bound}
 |W_{\mathbf t}^{\mathbf b}(n)|\leq W_{\mathbf t}(n)
 \leq L^{C(t_{\max})}.
\end{equation}
Indeed $q_0^{u-t}(u)_tV_g^{-t}$ is bounded uniformly in integers
$u\geq t$, with a bound depending only on $t_{\max}$ and the fixed
group data, and there are $O(T)$ groups.

Use the same choice at both endpoints in
\begin{equation}\label{sh:sign}
 \varepsilon(n)=1\quad\hbox{or}\quad
 \varepsilon(n)=(-1)^{\sum_{p\in\mathcal P}v_p(n)}.
\end{equation}
In either case $\varepsilon$ is completely multiplicative, real,
and of modulus one.  An invariant endpoint core has the form
\begin{equation}\label{sh:core}
 F(n)=\sum_{mpr=n_*}\alpha_m\chi_0(m)m^{i\sigma_0}
           \1_{p\in I_p,\ p\ {\rm prime}}p^{i\sigma_1}c_r.
\end{equation}
Fix $0<\tau<\eta<1/4$.  Here $I_p\subset[x^\tau,x^\eta]$ is an
interval, $|\alpha_m|,|c_r|\leq L^{C}$,
$\alpha_m=0$ unless $P^-(m)>W$, and the modulus of $\chi_0$ and
$|\sigma_0|,|\sigma_1|$ are at most $L^C$.  The constants in this
section may depend on these fixed bounds.  The word invariant refers
to the identity $F(vn)=F(n)$ whenever $v_*=1$.

\begin{theorem}[Twisted prime-slot correlation]\label{sh:correlation}
Let $F,G$ have the form \eqref{sh:core}, with possibly different
coefficients, characters, intervals, and frequencies, subject to fixed
bounds as above.  For every fixed $D_0>0$ there is a fixed $B>0$ with
the following property.  Suppose that the sequence $\alpha$ in $F$
satisfies
\begin{equation}\label{sh:discrepancy}
 \left|\sum_{m\in I}\alpha_m\chi(m)m^{-1+it}\right|\leq L^{-B}
 \quad\left(
 \begin{array}{l}
 I\subset[1,x^2]\text{ an interval},\\
 \operatorname{mod}(\chi)\leq L^B,\quad |t|\leq L^B
 \end{array}\right).
\end{equation}
There is no discrepancy requirement on $G$.  Uniformly for
$xL^{-C}\leq Z\leq xL^C$, integers $0<|k|\leq L^C$,
$|a_1|,|a_2|\leq L^C$, and smooth functions $\psi$ supported in a
fixed compact subinterval of $(0,\infty)$ with
$\|\psi^{(j)}\|_\infty\ll_j L^{C(j+1)}$, one has
\begin{equation}\label{sh:correlation-bound}
 \sum_{\substack{z>0\\z+k>0}}\frac{\psi(z/Z)}z
 z^{ia_1}(z+k)^{ia_2}\overline{F(z)}G(z+k)
 W_\ell(z)W_\ell(z+k)\varepsilon(z)\varepsilon(z+k)
 \ll L^{-D_0}.
\end{equation}
The exponent $B$ depends only on $D_0$ and the displayed fixed data.
\end{theorem}

The proof uses a residual operator made up of one raw term and signed
comparison terms. Transference makes the residual pairing small; after
the shared dilation is removed, the raw term is the correlation in
\eqref{sh:correlation-bound}, because the common sign cancels.
Each comparison has two independent large free products, allowing a
Fourier estimate. We will bound these comparisons and subtract them
from the residual pairing.

\subsection{The graph inputs}

We give the operator definitions to specify their normalizations.
For a sufficiently large fixed integer $J$, put $M=J+\ell$ and
$\mu_g(p)=1/(pV_g)$.  A pattern $\nu$ specifies
$C_g\subset\{1,\ldots,J\}$ with $C_g=\varnothing$ in small groups
and $|C_g|\leq m_0$ in big groups.  Put $t_g=\ell+|C_g|$.
The first $J$ slots outside $C_g$ are shared between the endpoints.
To form the dilation $D$, take their labels and independent auxiliary
labels of law $\mu_g$ in the slots of $C_g$.  Thus $D$ contains $J$
labels from each group, counted with multiplicity.
The coefficient $K_\nu$ may depend only on the ordered big-group
source and target labels and the auxiliary big-group free labels.

A physical state at $n\in\mathbb Z$ consists of $M$ ordered,
distinct divisors from each group.  Each state has mass
$\prod_gV_g^{-M}$, with counting measure in $n$.  For a row
transition from $n$ to $n'=n+kD$, retain the shared labels and sum
over all physical target lists with coefficient $\prod_gV_g^{-t_g}$.
Auxiliary free labels are forbidden from both endpoint lists; they
may coincide with one another.  The row multiplier is
\begin{equation}\label{sh:physical-multiplier}
 K_\nu D^{i\zeta}
 q_0^{(\omega_P(n)-Ms_P)/2}q_0^{(\omega_P(n')-Ms_P)/2}.
\end{equation}
The operator $\mathcal A_\zeta$ is the sum of these row operations,
averaged over independent permutations of the $M$ slots in every
group at each endpoint.  In particular the joint normalization of
the two lists in group $g$ is $V_g^{-M-t_g}$.

The ideal operator acts on the probability space of the ordered
big-group lists, with independent coordinates of law $\mu_g$;
repetitions are permitted.  It retains the same shared slots and
samples both the unshared target slots and the auxiliary slots
independently.  If $D_{\mathcal B}$ is the big-group part of $D$,
its multiplier is
$K_\nu D_{\mathcal B}^{i\zeta}\e(\Theta D_{\mathcal B})$.
The sum, symmetrized at both ends, is $\mathcal T_\zeta(\Theta)$.

\begin{proposition}[Transference input]\label{sh:transference}
The following are the Local Transference Theorem, its Endpoint Pairing
Corollary, and the Absolute Physical Bounds Lemma of
\cite[Section~3, Theorem~3.5, Corollary~3.11, and Lemma~3.4]{SmoothShifted}.
Assume $\sum_\nu\sup|K_\nu|\leq L^{C_1}$, where $m_0,C_1$ are
fixed independently of $J$.  For every fixed $E>0$ there is an
$E_{\rm id}>0$, depending only on $E$ and the group data, such that
\[
 \sup_{\Theta\in\mathbb R}
       \|\mathcal T_\zeta(\Theta)\|_{2\to2}\leq L^{-E_{\rm id}}
\]
implies the following conclusion for all sufficiently large fixed
$J$.  Let $I=[X,2X)$ with $x^\gamma\leq X\leq x^{1/\gamma}$ for
some fixed $\gamma>0$.  If $f$ is supported on positions in $I$,
is independent of the chosen marks, and
\[
 \sup|f|\leq\exp(O(\sqrt L)),\qquad
 \|f\|,\|g\|\leq X^{1/2}L^{C_3},
\]
then
\begin{equation}\label{sh:pairing}
 |\langle f,\mathcal A_\zeta g\rangle|
 \ll XL^{-E+2C_3}.
\end{equation}
The lower bound for $J$ may depend on $m_0,C_1$; the threshold for
$x$ may also depend on $J$, the fixed exponent bounding $|k|$, and
$\gamma$.  There is no additional restriction on $\zeta$ beyond the
ideal norm hypothesis.  The operator obtained by taking absolute
values of all transition coefficients has row and column sums at
most $L^{C_{\rm abs}}$, with $C_{\rm abs}$ independent of $J$.
\end{proposition}

\begin{proposition}[Residual ideal family and comparison kernels]
\label{sh:ideal}
The following are the Residual Ideal Operator Theorem and Comparison
Kernel Lemma of \cite[Section~4, Theorem~4.1 and Lemma~4.3]{SmoothShifted}.
For every prescribed $E_{\rm id}>0$ and fixed $C_4\geq0$, there are
fixed $m_0,J_0,C_1$ such that, for each fixed $J\geq J_0$, a family
consisting of the raw pattern $C_g=\varnothing$, $K_0=1$, and signed
comparison patterns satisfies
\begin{equation}\label{sh:ideal-bound}
 \sum_\nu\sup|K_\nu|\leq L^{C_1},\qquad
 \sup_{\substack{\Theta\in\mathbb R\\|\zeta|\leq L^{C_4}}}
       \|\mathcal T_\zeta(\Theta)\|_{2\to2}
       \leq L^{-E_{\rm id}}.
\end{equation}
The family is independent of the individual values of $\Theta,\zeta$,
and all its defining parameters are independent of $J$.

Split the big groups into two blocks.  Every comparison frees
nonempty sets of probes $A,B$ in the respective blocks and has
coefficient
\begin{equation}\label{sh:comparison-kernel-product}
 -(-1)^{|A|+|B|}\mathcal K_A(D_A,X_A)\mathcal K_B(D_B,X_B).
\end{equation}
Here $D_A,D_B$ are the free products, $X_A$ is the corresponding
target-mark product, and $X_B$ the corresponding source-mark product.
No shared label or final $\ell$ mark enters these kernels.  For a
slot set $\Lambda$, their exact form is
\begin{equation}\label{sh:kernel}
 \mathcal K_\Lambda(D,X)
 =\sum_{j:\nu_{\Lambda j}\geq\xi}
   \frac{\1_{\log D\in I_j}\1_{\log X\in I_j}}{\nu_{\Lambda j}}
   \sum_{\substack{\chi\ {\rm primitive}\\
                    \operatorname{cond}(\chi)\leq Q}}
               \overline{\chi(D)}\chi(X),
 \quad I_j=[jh,(j+1)h),
\end{equation}
where $\nu_{\Lambda j}$ is the probability of the indicated cell
under the independent slot laws.  The quantities $Q,h^{-1},\xi^{-1}$
are fixed powers of $L$, with $h\leq1$.  There are at most $Q^2$
characters and $O_{m_0}(1+TL^d/h)$ nonempty cells, and
\[
 \sup|\mathcal K_\Lambda|\leq Q^2\xi^{-1},\qquad
 \sup_X\mathbb E_D|\mathcal K_\Lambda(D,X)|\leq Q^2.
\]
More precisely, the parameters can be chosen, independently of $J$,
so that for any fixed $A_0>0$, any tuple tests $f(X_B),g(X_A)$,
and every $\theta\in\mathbb R$, $|\zeta|\leq L^{C_4}$,
\begin{align}
 \Big|\mathbb E\,\overline{f(X_B)}g(X_A)
 \big[&(X_AX_B)^{i\zeta}\e(\theta X_AX_B)\notag\\
 &-\mathcal K_A(D_A,X_A)\mathcal K_B(D_B,X_B)
       (D_AD_B)^{i\zeta}\e(\theta D_AD_B)\big]\Big|
 \leq L^{-A_0}\|f\|_2\|g\|_2.
 \label{sh:comparison-contract}
\end{align}
All four label tuples in this expectation are independent.  The tests
need not be functions only of their products.  Expanding the two
kernels in \eqref{sh:comparison-kernel-product}, including all
patterns, has total coefficient mass and number of terms bounded by
fixed powers of $L$.
\end{proposition}

These are precisely the graph inputs we use.  In particular,
\eqref{sh:kernel} consists of two global character and cell expansions,
not one expansion for each group.

\subsection{Endpoint norms and removal of the shared dilation}

We now prove Theorem~\ref{sh:correlation} using these contracts.
Write $\tau(v)$ for the divisor-counting function.  The pointwise
estimate
\begin{equation}\label{sh:core-divisor}
 |F(n)|\ll L^{C'}\tau(n_*)^2
\end{equation}
follows by counting the three-factor decompositions in
\eqref{sh:core}; the same estimate holds for $G$.
On physical states the function
\[
 F(n)\varepsilon(n)q_0^{(\omega_P(n)-Ms_P)/2}
\]
is independent of the chosen marks, and its damping is at most one.
For a fixed physical list of product $P$, one has $n_*\mid n/P$.
Consequently Proposition~\ref{pre:moments} gives, uniformly on a fixed
enlargement of a position dyad,
\[
 \sum_{\substack{n\asymp X\\P\mid n}}|F(n)|^2
 \ll L^{C'}\sum_{v\ll X/P}\tau(v)^4
 \ll (X/P)L^{C''}.
\]
The normalized reciprocal sum of the admissible lists is at most
\[
 \prod_g V_g^{-M}
       \sum_{\text{lists}}\frac1P
 \leq\prod_g\left(\sum_{p\in\mathcal P_g}\frac1{pV_g}\right)^M=1.
\]
Here $P\leq\exp(O_J(TL^d))=x^{o(1)}$, so $X/P$ is in the
ordinary divisor-moment range.  We have proved
\begin{equation}\label{sh:endpoint-norm}
 \|F\varepsilon q_0^{(\omega_P-Ms_P)/2}\|_{n\asymp X}
 \ll X^{1/2}L^{C_3},
\end{equation}
with $C_3$ independent of $J$ and of the subsequent kernel choices.
Bounded smooth factors do not affect this independence.  Moreover,
both $m$ and $p$ in \eqref{sh:core} are $W$-rough for large $x$.
There are at most $O(\sqrt L)$ prime factors above $W$ in
$n\asymp x^{1+o(1)}$, counted with multiplicity.  Assigning each such
factor to $m,p$, or the remaining factor shows
\begin{equation}\label{sh:endpoint-sup}
 |F(n)|\leq L^{C'}3^{O(\sqrt L)}=\exp(O(\sqrt L)).
\end{equation}
This proves every endpoint hypothesis of Proposition~\ref{sh:transference}.

Consider the physical pairing with extra cutoff, harmonic measure,
and phases
\begin{equation}\label{sh:lifted-pairing}
 \frac{\psi(n/(ZD))}{n}
 D^{-i(a_1+a_2)}n^{ia_1}(n')^{ia_2}
 \overline{F(n)}G(n')\varepsilon(n)\varepsilon(n')
 q_0^{(\omega_P(n)-Ms_P)/2}q_0^{(\omega_P(n')-Ms_P)/2}.
\end{equation}
The row operation itself contributes the other two half-damping
factors in \eqref{sh:physical-multiplier}.  Summation over states,
auxiliary draws, and patterns is understood in
\eqref{sh:lifted-pairing}, with the row coefficient $K_\nu$.

Here $n\asymp ZD=x^{1+o(1)}$ and $|n'-n|=x^{o(1)}$.  Split
$n$ among $O(L)$ dyads $[X,2X)$ and restrict $n'$ to a fixed
enlargement of each dyad.  As in Lemma~\ref{pre:separation}, Fourier inversion of
$u\mapsto\psi(e^u)$ separates the cutoff into powers of $n$ and
$D$; put the bounded factor $X/n$ in the first vector and divide
the pairing by $X$.  The total Fourier integral mass is a fixed
power of $L$, independent of $J$.

We specify a useful quantifier here.  If the logarithmic cutoff
derivatives are bounded by $O_j(L^{C_\psi(j+1)})$, fix an exponent
$F_0>C_\psi+2$ before forming the ideal family.  Integration by
parts gives, for each fixed $N$,
\[
 \int_{|v|>L^{F_0}}|\widehat{\psi\circ\exp}(v)|\,dv
 \ll_N L^{C_\psi(N+1)-F_0(N-1)}.
\]
Thus this same cutoff exponent permits arbitrary saving by increasing
$N$.  Choose $C_4$ to contain $a_1+a_2$ and all frequencies in the
truncated integral.  By \eqref{sh:endpoint-norm} and
\eqref{sh:endpoint-sup}, choose $E$ large enough for the desired
saving after the fixed Fourier and dyadic costs; then choose
$E_{\rm id}$, the residual family, and finally $J$.  The Fourier
tail is bounded using the absolute row and column bounds of
Proposition~\ref{sh:transference}; its accuracy can be increased
after the family and $J$ have been fixed without changing $F_0$ or
$C_4$.  It follows that the sum of the residual pairings is
$O(L^{-D})$ for any prescribed fixed $D$.

We next identify each individual pairing.  Write $D=D_RD_C$,
where $D_R$ is the shared product and $D_C$ the free product, and put
$n=D_Rw$, $n'=D_Rw'$.  Then $w'=w+kD_C$.  Away from overlaps of
shared labels with one another or with the divisors of $w,w'$, the
number of shared labels in group $g$ is $M-t_g$, and
\[
 \omega_g(D_Rw)-M=\omega_g(w)-t_g.
\]
The full damping therefore leaves exactly the two damping factors
of $W_{\mathbf t}(w)W_{\mathbf t}(w')$.  The identity
\[
 V_g^{-M-t_g}=V_g^{-(M-t_g)}V_g^{-2t_g}
\]
and $1/n=1/(D_Rw)$ turn the shared labels into independent draws
of law $\mu_g$, leaving the normalizations for the two unshared
marked lists.  The cores are unchanged.  The sign identity is exact,
even when overlaps occur:
\begin{equation}\label{sh:shared-sign}
 \varepsilon(D_Rw)\varepsilon(D_Rw')
 =\varepsilon(D_R)^2\varepsilon(w)\varepsilon(w')
 =\varepsilon(w)\varepsilon(w').
\end{equation}
The phases and cutoff in \eqref{sh:lifted-pairing} become
\[
 \psi(w/(ZD_C))\,(w/D_C)^{ia_1}(w'/D_C)^{ia_2}/(D_Rw).
\]
Since the outer vectors were mark-independent, the initial
symmetrizations do not alter this computation.  After summing the
shared labels, the pairing for pattern $\nu$ is, up to negligible
error,
\begin{align}
 \mathbb E_{\rm free}\sum_{\substack{w>0\\w'=w+kD_C>0}}
 &\frac{\psi(w/(ZD_C))}{w}
 (w/D_C)^{ia_1}(w'/D_C)^{ia_2}
 \overline{F(w)}G(w')\varepsilon(w)\varepsilon(w')\notag\\
 &\hspace{10mm}\cdot W_{\mathbf t}(w)W_{\mathbf t}(w')
                   \mathbb E_{\rm marks}K_\nu.
 \label{sh:stripped}
\end{align}
The final expectation means averaging over the unshared marked lists.

For completeness, all discarded restrictions have superpolynomial
logarithmic saving.  Conditional on $w,w'$, the free labels, and
earlier shared draws, at most $O_J(L)$ prime values are excluded
from a shared draw.  Every atom is $O(\exp(-L^a))$.  When bounding
actual overlapping configurations, unshared physical marks still
divide $w,w'$; ignoring damping changes their weight by at most
$q_0^{-2Ms_P}=L^{O_J(1)}$.  Equations \eqref{sh:weight-bound} and
\eqref{sh:core-divisor}, followed by Cauchy--Schwarz and fixed divisor
moments on the two translated intervals, bound the remaining harmonic
sum by $L^{O_J(1)}$.  Shared overlaps thus cost
$L^{O_J(1)}\exp(-L^a)$.

If an auxiliary free label $p'$ belongs to an unshared endpoint
list, then $p'\mid D_C$ and hence $p'\mid w,w'$.  Set $w=p'y$.
The new shift $kD_C/p'$ is integral, the invariant cores are
unchanged, and the harmonic measure supplies $1/p'$.  Applying the
same divisor moments at scale $ZD_C/p'$ bounds this contribution
by $L^{O_J(1)}/p'$.  There are $O_J(T)$ free slots and
$p'\geq\exp(L^c)$.  This bounds all free exclusions.  After summing
the fixed-power coefficient costs, the error in
\eqref{sh:stripped} is $O_A(L^{-A})$ for every fixed $A$.
For the raw pattern, $D_C=1$, $\mathbf t=(\ell,\ldots,\ell)$,
and \eqref{sh:stripped} is exactly the left side of
\eqref{sh:correlation-bound}.  It remains to bound the comparisons.

\subsection{Comparison multipliers and minor arcs}

Expand the two kernels in a comparison using \eqref{sh:kernel}.
Their four factors are bounded weights on $D_A,D_B$ and on the
big unshared marks at the respective endpoints, times a total
coefficient cost $L^{O(1)}$.  The free products are restricted to
logarithmic cells of width at most one.  Let their lower endpoints
be $U_1,U_2$, put $H'=U_1U_2$, and set $Y=ZH'$.  Then
\begin{equation}\label{sh:comparison-scales}
 U_1,U_2\geq\exp(L^c-O(1)),\quad
 H'\leq\exp(O(TL^d)),\quad Y=x^{1+o(1)}.
\end{equation}
The positions $w,w'$ in \eqref{sh:stripped} are comparable to $Y$.
Logarithmic Fourier separation of the smooth cutoff and phases
therefore reduces it, with total cost $L^{O(1)}/Y$, to expressions
\begin{equation}\label{sh:comparison-convolution}
 \mathbb E\,b_1(D_A)b_2(D_B)
       \sum_w\overline{\mathcal U(w)}\mathcal V(w+kD_AD_B),
 \qquad |b_1|,|b_2|\leq1.
\end{equation}
An endpoint here is its core times $\varepsilon$, a power twist of
fixed log-power frequency, $W_{\mathbf t}^{\mathbf b}$ with
$\mathbf b$ depending only on big marks, and a smooth size cutoff
at $Y$.  Fixed divisor moments give
\begin{equation}\label{sh:comparison-norms}
 \sum_w|\mathcal U(w)|^2+\sum_w|\mathcal V(w)|^2
 \ll YL^{C_5}.
\end{equation}
They also control the Fourier-separation tails to arbitrary accuracy.

The Fourier multiplier of \eqref{sh:comparison-convolution} is
\[
 \mathfrak m(\theta)=\mathbb E\,b_1(D_A)b_2(D_B)
                                     \e(k\theta D_AD_B),
 \qquad |\mathfrak m(\theta)|\leq1.
\]
For a selected product $D$ with at most $m_0$ slots per group,
unique factorization gives
\begin{equation}\label{sh:product-atoms}
 \mathbb P(D=n)=\frac1n\prod_g
       \frac{k_g!}{V_g^{k_g}\prod_{p\in\mathcal P_g}v_p(n)!}
 \leq\frac{L^{C_6}}n.
\end{equation}
Thus a bounded test collapsed onto product values $n\asymp U_i$
has squared coefficient norm $O(L^{C_6}/U_i)$.

We use the integer bilinear estimate of
\cite[Section~4, Lemma~4.2]{SmoothShifted}: for coefficients on
intervals of lengths $O(U),O(V)$, and
$\alpha=u/r+\beta$ with $(u,r)=1$, $|\beta|\leq r^{-2}$,
\[
 \left|\sum_{n,t}a_nb_t\e(\alpha nt)\right|
 \ll\|a\|_2\|b\|_2
       [U+(1+V/r)\{U+r\log(2r)\}]^{1/2}.
\]
Arbitrary missing coefficients and translated containing intervals
are allowed.  Together with \eqref{sh:product-atoms}, this bounds
$|\mathfrak m(\theta)|$ by
\begin{equation}\label{sh:minor-bound}
 L^{C_7}\left(\frac1{U_2}+\frac1r+
          \frac{\log(2r)}{U_1}+
          \frac{r\log(2r)}{H'}\right)^{1/2}
\end{equation}
whenever $r$ is a Dirichlet-approximation denominator for $k\theta$.

Choose such an approximation with maximum denominator
$\lfloor H'/L^{C'}\rfloor$.  If $r>L^{C'}$, all terms in
\eqref{sh:minor-bound} save an arbitrarily large fixed log power
when $C'$ is sufficiently large, by \eqref{sh:comparison-scales}.
If $r\leq L^{C'}$, then $\theta$ belongs, after accounting for
$|k|\leq L^C$, to an arc
\begin{equation}\label{sh:arcs}
 \left|\theta-\frac hr\right|_{\mathbb R/\mathbb Z}
 \leq\frac{L^{C_{\rm arc}}}{H'},
 \qquad 1\leq r\leq L^{C_{\rm arc}},\quad h\bmod r,
\end{equation}
for a sufficiently large fixed $C_{\rm arc}$.  The letters $h,r$
in \eqref{sh:arcs} denote the resulting rational center, rather
than necessarily the original approximation.  There are a fixed
logarithmic power of arcs.  Parseval and
\eqref{sh:comparison-norms} dispose of their complement.  On the
arcs, Cauchy--Schwarz shows that it suffices to prove, for every
fixed $A>0$ and every arc center $h/r$,
\begin{equation}\label{sh:arc-target}
 \int_{|\beta|\leq1/H}
   |\widehat{\mathcal U}(h/r+\beta)|^2\,d\beta
 \ll_A YL^{-A},\qquad H=H'/L^{C_{\rm arc}}.
\end{equation}
Here $\widehat{\mathcal U}(\theta)=\sum_w\mathcal U(w)\e(\theta w)$,
$H\to\infty$, and $H=x^{o(1)}$.

\subsection{The endpoint on a major arc}

Choose a small group $g_s$ and a designated one of its
$\ell\geq1$ marks.  Because the mark weight $\mathbf b$ uses only
big marks, removing this prime gives, unless a prime in $g_s$
divides $w$ twice,
\begin{equation}\label{sh:small-mark-extraction}
 W_{\mathbf t}^{\mathbf b}(w)
 =\frac1{V_{g_s}}\sum_{\substack{p_s\mid w\\p_s\in\mathcal P_{g_s}}}
           W_{\mathbf t-\mathbf e_{g_s}}^{\mathbf b}(w/p_s).
\end{equation}
The equality includes damping: removing $p_s$ decreases both
$\omega_P$ and $\sum_g t_g$ by one.  On the excluded set both
weights are bounded by fixed log powers.  Since $w_*\mid w/p_s^2$
on a multiple of $p_s^2$, divisor moments show that the endpoint
error has squared norm
\begin{equation}\label{sh:small-square-error}
 \ll YL^{C_8}\sum_{p_s\in\mathcal P_{g_s}}p_s^{-2}
 \ll YL^{C_8}\exp(-L^a).
\end{equation}
Its Fourier energy is negligible by Parseval.

The factors $m,p$ in \eqref{sh:core} contain no group prime.
After \eqref{sh:small-mark-extraction}, write $w=mp_sp r'$.
Complete multiplicativity gives
$\varepsilon(w)=\varepsilon(p_s)\varepsilon(r')$; the remaining
marked weight and $c_{r'_*}$ are functions of $r'$ alone.  They
have a fixed log-power bound, so may be absorbed into one residual
coefficient.

To remove the rational center, fix $d'=(r',r)$ and put $q=r/d'$.
The first three factors $m,p_s,p$ are units modulo $r$ for all
large $x$, and $r'/d'$ is a unit modulo $q$.  On units the exact
character expansion is
\begin{equation}\label{sh:rational-characters}
 \e(hv/q)=\sum_{\chi\bmod q}c_{h,q}(\chi)\chi(v),\qquad
 c_{h,q}(\chi)=\frac1{\varphi(q)}
       \sum_{v\bmod q}^{*}\e(hv/q)\overline{\chi(v)},
 \quad |c_{h,q}(\chi)|\leq1.
\end{equation}
Apply this with $v=mp_sp(r'/d')$.  The gcd restriction and
$\chi(r'/d')$ go into the residual coefficient.  Summing over
$d'$ and all characters costs a fixed power of $L$ and includes
the principal and imprimitive characters.  The product of the new
character on $m$ with $\chi_0$ is a character of fixed log-power
modulus.  It is therefore enough to bound the energy at zero of
\begin{equation}\label{sh:four-factors}
 a_w=\psi_1(w/Y)w^{iv_1}
       \sum_{mp_sp r'=w}\alpha_m\chi_m(m)m^{i\sigma_0}
       b_s(p_s)\1_{p\in I_p,\ p\ {\rm prime}}
       \chi_p(p)p^{i\sigma_1}d_{r'}.
\end{equation}
Here $p_s\in\mathcal P_{g_s}$, $|b_s|\ll1$,
$|d_{r'}|\leq L^{C_9}$, and all characters, frequencies and
smooth-cutoff derivatives have fixed log-power bounds.  The sign
on $p_s$ and $1/V_{g_s}$ are included in $b_s$.  In particular,
\begin{equation}\label{sh:a-norm}
 \sum_w|a_w|^2\ll YL^{C_{10}}.
\end{equation}

\subsection{Local Fourier energy and Mellin polynomials}

We include the scale conversion from
\cite[Section~5, equation~(5.16)]{SmoothShifted} with its proof,
as its normalization is useful here.

\begin{lemma}[Local Mellin energy]\label{sh:mellin-energy}
Let $a_w$ be supported on $w\asymp Y$, where $Y=x^{1+o(1)}$,
and let $H=x^{o(1)}$ tend to infinity.  For every fixed $j\geq1$,
\begin{equation}\label{sh:mellin-inequality}
 \int_{|\beta|\leq1/H}|\widehat a(\beta)|^2\,d\beta
 \ll_j \frac1Y\int_{\mathbb R}(1+H|t|/Y)^{-j}
           \left|\sum_w a_ww^{it}\right|^2dt
       +\left(\frac HY\right)^2\sum_w|a_w|^2.
\end{equation}
\end{lemma}

\begin{proof}
Take a smooth bump $K$ of integral one, supported sufficiently near
zero that its Fourier transform, with convention
$\widehat K(\xi)=\int K(u)\e(-\xi u)\,du$, has modulus at least
$1/2$ on $[-1,1]$.  Plancherel gives
\[
 \int_{|\beta|\leq1/H}|\widehat a(\beta)|^2\,d\beta
 \ll H^{-2}\int_{\mathbb R}
          \left|\sum_w a_wK((w-v)/H)\right|^2dv.
\]
Only $v\asymp Y$ contributes.  Keep this restriction when replacing
the kernel by $K(v\log(w/v)/H)$.  On the union of their supports,
$|w-v|\ll H$, and their arguments differ by $O(H/Y)$.
Cauchy--Schwarz over the $O(H)$ possible integers $w$, followed
by integration over the $O(H)$ centers for each $w$, bounds the
normalized squared error by $(H/Y)^2\sum_w|a_w|^2$.

Put $h_0=H/Y$, $v=Y\exp(u)$, and $P_a(t)=\sum_w a_ww^{it}$.
Fourier inversion for the new kernel gives
\[
 \sum_w a_wK(v\log(w/v)/H)
 =\frac{h_0}{2\pi}\int_{\mathbb R}
   e^{-u}\widehat K(h_0e^{-u}t/(2\pi))v^{-it}P_a(t)\,dt.
\]
Insert a fixed smooth compact cutoff in $u$ equal to one on the
required range.  The Fourier transform in $u$ of the resulting
amplitude $e^{-u}\widehat K(h_0e^{-u}t/(2\pi))$ is bounded by
\[
 C_j(1+|s|)^{-2}(1+h_0|t|)^{-j},
\]
by two integrations by parts and the Schwartz bounds for
$\widehat K$.  Expand that amplitude in its $u$-Fourier transform,
apply Minkowski in $s$ and Plancherel in $u$, and use $dv\ll Y\,du$.
The normalization is $H^{-2}Yh_0^2=Y^{-1}$.  Increasing the decay
order if necessary proves \eqref{sh:mellin-inequality}.
\end{proof}

Apply this lemma to \eqref{sh:four-factors}.  The error term is
$O_A(YL^{-A})$ for every fixed $A$, by \eqref{sh:a-norm}.
Dirichlet-polynomial mean squares imply, on every dyadic time scale
$R\geq1$,
\begin{equation}\label{sh:mellin-tail-ms}
 Y^{-2}\int_{R\leq|t|\leq2R}|P_a(t)|^2dt
 \ll(1+R/Y)L^{C_{11}}.
\end{equation}
Consequently, with $T_0=LY/H$, the part $|t|>T_0$ of the first
term in \eqref{sh:mellin-inequality} is at most
\[
 YL^{C_{11}}\sum_{r\geq0}
       (2^rL)^{-j}\left(1+\frac{2^rL}{H}\right)
 \ll_j YL^{C_{11}-j}(1+L/H).
\]
It is $O_A(YL^{-A})$ on taking $j$ large enough.

Split the four variables in \eqref{sh:four-factors} into dyads.
There are $O(L^4)$ relevant boxes, and their product scales are
comparable to $Y$.  Write
\[
 \frac{P_a(t)}Y
 =\sum_w\frac1w\big((w/Y)\psi_1(w/Y)\big)w^{i(t+v_1)}
        \sum_{mp_sp r'=w}(\text{the four coefficients}).
\]
Fourier-separate the parenthesized function in $\log(w/Y)$.  On
every box the fourfold collapsed reciprocal coefficients have
squared norm $\ll L^{C_{12}}/Y$, by a fixed divisor moment.
Hence centered mean squares give a bound $L^{C_{13}}$ for the
integral of their squared polynomial on any translate of
$[-T_0,T_0]$.  Minkowski therefore bounds a separating-frequency
tail by a fixed log power times the square of the $L^1$ mass of
that tail.  Integration by parts makes it arbitrarily small at a
fixed log-power cutoff.  This proves the claimed separation also
for unbounded tail frequencies.  Retained shifts enlarge the
time interval to at most $[-2T_0,2T_0]$, since $Y/H$ exceeds every
fixed log power.

After absorbing these shifts into $t$ and the fixed frequencies,
it suffices to prove, for arbitrary fixed $A'>0$,
\begin{equation}\label{sh:polynomial-target}
 \int_{|t|\leq2LY/H}|P_s(t)P_l(t)\mathcal R(t)|^2dt
 \ll L^{-A'},
\end{equation}
where
\begin{align*}
 P_s(t)&=\sum_{p_s\asymp P}b_s(p_s)p_s^{-1+it},\\
 P_l(t)&=\sum_{\substack{p\asymp P',\ p\in I_p\\p\ {\rm prime}}}
                      \chi_p(p)p^{-1+i\sigma_1+it},\\
 \mathcal R(t)&=
  \left(\sum_{m\asymp M_1}\alpha_m\chi_m(m)m^{-1+i\sigma_0+it}\right)
  \left(\sum_{r'\asymp R_1}d_{r'}(r')^{-1+it}\right).
\end{align*}
The small primes retain their group restriction, and each dyad is
intersected with its original interval.  Their sizes satisfy
\begin{equation}\label{sh:polynomial-scales}
 \tfrac12\exp(L^a)\leq P\leq\exp(L^b),\qquad
 x^\tau/2\leq P'\leq x^\eta,\qquad PM_1P'R_1\asymp Y.
\end{equation}
For any subproduct of these four polynomials, collapsed as
$\sum_{n\asymp U}c_nn^{it}$, a fixed divisor moment gives
\begin{equation}\label{sh:collapsed-norm}
 \sum_n|c_n|^2\ll L^{C_{14}}/U.
\end{equation}
Only four convolution factors are involved: all remaining marked
weights have already entered the single bounded coefficient $d_{r'}$.

\subsection{Exceptional times and completion of the proof}

First consider $|P_s(t)|\leq L^{-A_s}$.  The polynomial
$P_l\mathcal R$ has size $Y/P$, and
\[
 \frac{Y/P}{2LY/H}=\frac{H}{2LP}\longrightarrow\infty
\]
by $b<c$ and \eqref{sh:comparison-scales}.  Its mean square on
the entire time interval is $L^{O(1)}$, by
\eqref{sh:collapsed-norm}.  Choose $A_s$ sufficiently large to
bound this part of \eqref{sh:polynomial-target}.

Let $\mathcal J$ be the integer unit intervals meeting
$\{|t|\leq2LY/H:|P_s(t)|>L^{-A_s}\}$.  We establish
\begin{equation}\label{sh:exceptional-count}
 |\mathcal J|\leq
 \exp\big(O(L^b+L^{1-a}\log L)\big)=Y^{o(1)}.
\end{equation}
Put $j_0=\lfloor\log Y/\log(2P)\rfloor$.  Unique factorization
shows directly that the squared coefficient norm of $P_s^{j_0}$
is at most
\begin{equation}\label{sh:factorial-norm}
 j_0!\left(\sum_{p_s\asymp P}|b_s(p_s)|^2p_s^{-2}\right)^{j_0}
 \leq j_0!(C/P)^{j_0}.
\end{equation}
Indeed for each multiset of $j_0$ primes, one of the two multinomial
coefficients in its squared coefficient is at most $j_0!$.
This argument remains valid when $j_0$ grows; no growing-order
divisor-moment estimate is being used.

Select a point with $|P_s|>L^{-A_s}$ in each occupied interval.
Split the intervals into three classes according to their integer
index modulo three, so the chosen points in a class are separated
by at least one.  The indices of $P_s^{j_0}$ are at most $Y$, and
the time range is contained in $[-Y,Y]$ for large $x$.  The
separated-point mean-square inequality and \eqref{sh:factorial-norm}
give
\[
 |\mathcal J|\ll YL^{C_{15}}j_0!
                         (CL^{2A_s}/P)^{j_0}.
\]
Since $YP^{-j_0}\leq2P\,2^{j_0}$,
$\log P\leq L^b$, and $j_0=O(L^{1-a})$, this proves
\eqref{sh:exceptional-count}.

We next prove the sparse mean-square estimate
\begin{equation}\label{sh:sparse-supremum}
 \sum_{I\in\mathcal J}\sup_{t\in I}|\mathcal R(t)|^2
 \ll L^{C_{16}}.
\end{equation}
Collapse $\mathcal R(t)=\sum_{n\asymp U}c_nn^{it}$, where
\[
 U=M_1R_1\asymp Y/(PP')\geq Y^{1-\eta-o(1)}>Y^{0.6},
 \qquad \sum_n|c_n|^2\ll L^{C_{14}}/U.
\]
Choose a maximizing point on each closed unit interval and again
split into three separated classes.  The Gram matrix of the
evaluation vectors $(n^{it_i})_{n\asymp U}$ has entries
$\sum_{n\asymp U}n^{i(t_i-t_j)}$.  For
$1\leq|\Delta|\leq c_*U$, with a sufficiently small fixed $c_*>0$,
the derivative of $\Delta\log n/(2\pi)$ is monotone, has magnitude
comparable to $|\Delta|/U$, and is bounded away from nonzero
integers.  The first derivative test of
Proposition~\ref{pre:derivatives} gives $O(U/|\Delta|)$.
For $c_*U\leq|\Delta|\ll Y$, the second derivative test gives
\[
 O\big(\sqrt{|\Delta|}+U/\sqrt{|\Delta|}\big)=O(Y^{1/2}).
\]
The diagonal entry is $O(U)$.  Thus all entries are bounded by
\[
 O\left(\frac{U}{1+|t_i-t_j|}+Y^{1/2}\right).
\]
Separation and \eqref{sh:exceptional-count} bound an absolute row
sum by
\[
 O\big(U\log Y+|\mathcal J|Y^{1/2}\big)=O(U\log Y).
\]
The operator norm is bounded by this row sum.  Multiplying by
$\sum|c_n|^2$ proves \eqref{sh:sparse-supremum}.

On the exceptional intervals, $|P_s(t)|\ll1$.  Prescribe a
sufficiently large saving for $P_l$, and apply
Lemma~\ref{lp:prime-polynomial}.  Choose $B_1$ larger than its
frequency threshold and the fixed exponent bounding $|\sigma_1|$.
For $|t|\geq L^{B_1}$ in the present range, the lemma applies to
$t+\sigma_1$: it is above the required logarithmic threshold and
\[
 |t+\sigma_1|\leq2LY/H+|\sigma_1|<x^2.
\]
It bounds $P_l$ by an arbitrarily large negative log power.
Equation \eqref{sh:sparse-supremum} then proves the required
integral bound for these times.

For $|t|\leq L^{B_1}$, use \eqref{sh:discrepancy} on the
$m$-polynomial in $\mathcal R$.  Its dyad lies below $x^2$;
the character modulus and shifted frequency are fixed log powers.
All other factors have fixed log-power absolute bounds.  Choosing
$B$ last, larger than these modulus and frequency exponents and
with $2B>B_1+C_{17}+A'$, makes this last integral $O(L^{-A'})$.
This proves \eqref{sh:polynomial-target}, hence
\eqref{sh:arc-target}.

The major-arc estimate, the minor-arc estimate, and Parseval now
bound every comparison \eqref{sh:stripped} by an arbitrarily large
negative log power, after paying the fixed number of arcs, kernel
expansions, and Fourier integrals.  Their sum has the same bound.
Subtracting it from the controlled residual pairing leaves the raw
pattern, which is \eqref{sh:correlation-bound}.  This completes
the proof of Theorem~\ref{sh:correlation}.

\begin{remark}[Order of choices]\label{sh:choice-order}
The endpoint norm exponent and the original Fourier cutoff exponent
are fixed from the input data.  Then choose the physical accuracy
$E$, ideal accuracy $E_{\rm id}$, ideal frequency range $C_4$, all
kernel parameters, and finally a sufficiently large fixed $J$.
Once this family is fixed, choose the minor-arc and major-arc
accuracies, the prime-polynomial accuracy and threshold $B_1$, and
last the discrepancy exponent $B$.  Any stronger Fourier-tail
saving required by these choices is obtained by more integrations
by parts at the already fixed cutoff exponent.  No kernel parameter
or frequency range has to be changed after choosing $J$.
\end{remark}
